\documentclass[11pt]{article}
\usepackage[margin=1in]{geometry}
\usepackage{amsmath,amssymb,amsthm}
\usepackage[hidelinks]{hyperref}
\newtheorem{proposition}{Proposition}
\newtheorem{corollary}{Corollary}
\newcommand{\df}{\mathrm{df}}
\newcommand{\tr}{\operatorname{tr}}
\newcommand{\Eg}{E_{\mathrm{gen,DE}}}
\title{Generalization at Accentuation-Optimal Regularization:\\A Conditional Effective-Dimension Universality}
\author{}
\date{September 18, 2026}
\begin{document}
\maketitle

\paragraph{Scope.}
This note concerns pixel-space ridge and the leading deterministic-equivalent
(DE) formulas. The identities below are exact identities between these DE
expressions, not exact finite-sample risk identities. ``Universality'' means
that teacher dependence can be absorbed into the calibration-selected effective
dimension; it does not assert a new distributional universality theorem.

\section{Setup and the calibration constraint}
Let $\Sigma\succeq0$ have eigenvalues $s_j$, and let $\beta_j^*$ be the
coordinates of a fixed teacher in its eigenbasis. Assume a finite spectrum,
$S=\sum_j s_j\beta_j^{*2}>0$, and response-noise variance $\sigma^2>0$.
The ridge objective is $n^{-1}\|y-X\beta\|^2+\lambda\|\beta\|^2$.
Define
\[
 B_{p,q}(\kappa)=\sum_j\frac{\beta_j^{*2}s_j^p}{(s_j+\kappa)^q},\qquad
 \df_{p,q}(\kappa)=\sum_j\frac{s_j^p}{(s_j+\kappa)^q},\qquad
 \df_1=\df_{1,1},\quad \df_2=\df_{2,2}.
\]
On the physical branch,
\[
 \lambda(\kappa)=\kappa\left(1-\frac{\df_1(\kappa)}n\right),\qquad
 \kappa>\kappa_0,\quad n-\df_1>0,\quad n-\df_2>0.
\]
Here $\kappa_0=0$ if $\operatorname{rank}\Sigma\le n$; otherwise it is the
positive solution of $\df_1(\kappa_0)=n$.
Write $V=\Eg+\sigma^2$. The generalization and norm/overlap DEs are
\begin{align}
 \Eg&=\frac{n(\kappa^2B_{1,2}+\sigma^2)}{n-\df_2}-\sigma^2,
 &\mu_N&=B_{1,1},
 &\mu_D&=B_{2,2}+\frac{\df_{1,2}}nV.
\end{align}
Since $B_{1,1}-B_{2,2}=\kappa B_{1,2}$, the condition
$z=\mu_D/\mu_N-1=0$ is equivalent to
\begin{equation}\label{eq:calibration}
 V=n\kappa\frac{B_{1,2}}{\df_{1,2}},\qquad
 \sigma^2=n\lambda(\kappa)\frac{B_{1,2}}{\df_{1,2}}.
\end{equation}
At such a point the leading quantities satisfy
$E_{\mathrm{acc}}^{(0)}/S=(z/(1+z))^2=0$,
$(R_{\mathrm{acc}}^2)^{(0)}=1-z^2=1$, and
$\mathrm{slope}_{\mathrm{acc}}^{(0)}=1/(1+z)=1$.

For completeness, a zero exists: the right-hand side of the second equation
in \eqref{eq:calibration} is continuous, tends to zero as
$\kappa\downarrow\kappa_0$, and tends to infinity as $\kappa\to\infty$,
because $B_{1,2}/\df_{1,2}\to S/\tr\Sigma>0$ and $\lambda\sim\kappa$.
The zero need not be unique. This is an existence statement for the DE
expressions under consideration, not a finite-sample calibration guarantee.

\section{The effective-dimension risk law}
\begin{proposition}[Generalization at an accentuation zero]
At any calibration point $\kappa_a>\kappa_0$ with $\lambda_a=\lambda(\kappa_a)$,
\begin{equation}\label{eq:main}
 \boxed{\quad
 \Eg(\kappa_a)=\sigma^2\left(\frac{\kappa_a}{\lambda_a}-1\right)
 =\sigma^2\frac{\df_1(\kappa_a)}{n-\df_1(\kappa_a)}.
 \quad}
\end{equation}
\end{proposition}
\begin{proof}
Divide the two equations in \eqref{eq:calibration} to obtain
$V/\sigma^2=\kappa_a/\lambda_a$, and then use the physical-branch relation.
\end{proof}

\paragraph{A useful off-calibration identity.}
For any $\lambda>0$ on the branch, direct algebra gives
\[
 \frac{\Eg+\sigma^2}{\sigma^2}-\frac{\kappa}{\lambda}
 =-\frac{\kappa^2B_{1,1}}{\lambda\sigma^2}\,z.
\]
Thus the risk law is equivalent to calibration, not a formula valid at all
regularizations. In particular it need not hold at a prediction optimum.

\section{Why the result is structurally special}
\subsection{From squared shrinkage to effective degrees of freedom}
Set $q_j=s_j/(s_j+\kappa)$. Then
\[
 \df_1=\sum_jq_j,\qquad \df_2=\sum_jq_j^2,\qquad
 \df_1-\df_2=\sum_jq_j(1-q_j)=\kappa\df_{1,2}.
\]
The general risk separates into a signal-dependent contribution and a
response-noise contribution:
\begin{equation}\label{eq:risk-decomp}
 \Eg=\frac{n\kappa^2B_{1,2}}{n-\df_2}
       +\frac{\sigma^2\df_2}{n-\df_2}.
\end{equation}
The first term includes the effect of random design on the signal; it should
not be called solely the squared bias of the mean estimator.
At calibration, the teacher-dependent shrinkage term obeys
\begin{equation}\label{eq:bias}
 \boxed{\kappa_a^2B_{1,2}(\kappa_a)
 =\sigma^2\frac{\df_1(\kappa_a)-\df_2(\kappa_a)}
 {n-\df_1(\kappa_a)}.}
\end{equation}
Substitution into \eqref{eq:risk-decomp} yields \eqref{eq:main}.
The nonnegative signal contribution effectively promotes the complexity
parameter from $\df_2$ to $\df_1$. Bias and variance do not cancel in the
prediction MSE. Instead, calibration fixes their relation so that their sum
has a single effective-dimension expression.

The difference $\sum_jq_j(1-q_j)$ is supported on partially retained modes;
it vanishes for a hard keep/discard filter. It quantifies the gap between
linear and squared shrinkage. At calibration, the teacher's shrinkage loss
is tied precisely to this gap and to the noise variance.

\subsection{An error self-consistency relation}
The general risk equation is equivalently
\[
 \Eg=\kappa^2B_{1,2}+\frac{\df_2}{n}V.
\]
Calibration substitutes
$\kappa^2B_{1,2}=(\df_1-\df_2)V/n$, hence
\begin{equation}
 \Eg=\delta_a V,\qquad V=\sigma^2+\delta_a V,
 \qquad \delta_a=\frac{\df_1(\kappa_a)}n\in(0,1).
\end{equation}
Consequently $V=\sigma^2/(1-\delta_a)$. This is an algebraic feedback
interpretation of the DE, not a literal iterative training mechanism.
It explains why perfect leading accentuation calibration can coexist with
large prediction error when effective complexity approaches the sample budget.

\subsection{OLS-like, but not an OLS equivalence theorem}
Equation \eqref{eq:main} has the familiar leading-order OLS form
$\sigma^2p_{\mathrm{eff}}/(n-p_{\mathrm{eff}})$ with
$p_{\mathrm{eff}}=\df_1(\kappa_a)$. This does not identify ridge with a
hard-truncated OLS estimator, assert unbiasedness, or give the exact
finite-sample OLS expectation (whose finite-size denominators can differ).
It is a risk-level structural analogy: shrinkage and estimation effects
combine into an effective-complexity law at this particular constraint.

\section{Weight-space geometry: calibration is radial optimality}
For a fixed nonzero estimated weight, let
$D=\|\hat\beta\|^2$ and $N=\beta^{*\top}\hat\beta$. Then
\[
 D=N\quad\Longleftrightarrow\quad
 \hat\beta^\top(\beta^*-\hat\beta)=0.
\]
Moreover, for $L(c)=\|c\hat\beta-\beta^*\|^2$,
\[
 L'(1)=2(D-N),\qquad \arg\min_c L(c)=N/D.
\]
Thus exact calibration means that the fitted weight already has the optimal
global Euclidean scale along its own direction. It does not imply optimal
direction or teacher recovery. The analogous expectation-level statement is
$\arg\min_c\mathbb E L(c)=\mathbb E N/\mathbb E D$ when the moments exist;
the leading DE substitutes $\mu_N,\mu_D$ for these quantities.

By contrast, prediction-optimal rescaling of a fixed weight uses
$\beta^{*\top}\Sigma\hat\beta/(\hat\beta^\top\Sigma\hat\beta)$.
The distinction is Euclidean versus data-covariance geometry. Even if
$\beta^*=\hat\beta+u$ with $u\perp\hat\beta$, prediction error is
$u^\top\Sigma u$, which need not be small.

\section{What is universal, and what remains teacher-dependent?}
At fixed spectrum, noise, sample size, and a shared calibration point
$\kappa_a$, all teachers satisfying \eqref{eq:calibration} have the same
DE generalization error. The constraint fixes the relevant quadratic teacher
statistic at that point. Changing the teacher generally changes the location
of the calibration point, and hence $\df_1(\kappa_a)$.

Accordingly the result is a \emph{teacher-free risk expression at a
teacher-dependent operating point}, not a teacher-independent numerical risk.
Calibration is an oracle condition unless its needed statistics can be
estimated. Finding a zero does not by itself supply a practical CV procedure.

\section{Constructive oracle bound and simultaneous performance}
\begin{corollary}
Let $E_{\mathrm{gen}}^\star=\inf_{\lambda\ge0}\Eg(\kappa(\lambda))$.
Every calibration point supplies a feasible witness, so
\[
 0\le E_{\mathrm{gen}}^\star\le
 \min\left\{S,\ \sigma^2\frac{\df_1(\kappa_a)}{n-\df_1(\kappa_a)}\right\}.
\]
The $S$ bound follows from the zero-predictor limit $\lambda\to\infty$.
\end{corollary}
The optimization inequality itself is elementary; what is constructive is
the existence of a calibration solution with a computable risk and a second,
simultaneously satisfied performance objective.
Writing $\rho=\sigma^2/S$, one obtains
\[
 \frac{E_{\mathrm{gen}}^\star}{S}
 \le\min\{1,\rho\delta_a/(1-\delta_a)\},\qquad
 \frac{\Eg(\kappa_a)}S\le\varepsilon
 \Longleftrightarrow
 \delta_a\le\frac{\varepsilon}{\rho+\varepsilon}.
\]
Thus simultaneous perfect leading accentuation and small prediction error
require sufficiently small effective complexity relative to the noise level.

If two distinct calibration zeros satisfy $\kappa_{a,2}>\kappa_{a,1}$,
then $\df_1'(\kappa)=-\df_{1,2}<0$ implies
$\Eg(\kappa_{a,2})<\Eg(\kappa_{a,1})$.
This orders risks \emph{among calibration zeros}, not along the whole ridge
path. The bound concerns the continuous DE oracle; it is not a per-trial
guarantee for empirical RidgeCV or an arbitrary finite selection grid.

\section{Limits of the statement}
\begin{itemize}
\item Finite-sample ratio fluctuations remain. With $R=N/D$,
$\mathbb E[(1-R)^2]=(1-\mathbb ER)^2+\operatorname{Var}(R)$.
$\mu_N=\mu_D$ does not force this risk to vanish. A distributional optimum
can move away from the leading calibration point.
\item The displayed ratios assume $\sigma^2>0$ and $\lambda_a>0$.
Zero noise requires boundary limits, especially when $\operatorname{rank}\Sigma=n$.
In overparameterized problems, $\df_1(\kappa_a)\to n$ can offset a vanishing
noise factor, so the formula does not imply noiseless recovery.
\item No monotonicity of the oracle gap or of calibration-selected effective
dimension is asserted without additional spectral/alignment assumptions.
\item General feature regression uses the input-space metric $FF^\top$ in
the norm term. The pixel-space cancellations here do not automatically extend
to a nonorthogonal feature map.
\end{itemize}

\paragraph{Takeaway.}
At leading order, zero accentuation error imposes a calibration constraint
that converts the teacher-dependent shrinkage contribution into an
effective-complexity contribution. Prediction risk then follows the OLS-like
law $\sigma^2p_{\mathrm{eff}}/(n-p_{\mathrm{eff}})$, with teacher dependence
retained through the calibration-selected effective dimension.
\end{document}
